%\ifdefined\handout
%\documentclass[handout,aspectratio=1610,xcolor={usenames,dvipsnames,table}]{beamer}
%\else
\documentclass[aspectratio=1610,xcolor={usenames,dvipsnames,table}]{beamer}
%\fi

\newcommand{\mylecture}{Configuration Management}

%\ifdefined\handout
%\documentclass[handout,aspectratio=1610,xcolor={usenames,dvipsnames,table}]{beamer}
%\else
\documentclass[aspectratio=1610,xcolor={usenames,dvipsnames,table}]{beamer}
%\fi

\newcommand{\mylecture}{Configuration Management}

%\ifdefined\handout
%\documentclass[handout,aspectratio=1610,xcolor={usenames,dvipsnames,table}]{beamer}
%\else
\documentclass[aspectratio=1610,xcolor={usenames,dvipsnames,table}]{beamer}
%\fi

\newcommand{\mylecture}{Configuration Management}

\input{../setup}
\input{../shared/setup}

\lstdefinelanguage{dump}
{
	morekeywords={kdbOpen,ksNew,keyNew,keyMeta,keyCopyMeta,keyEnd,ksEnd,kdbClose},
	sensitive=false,
	morecomment=[l]{//},
	morecomment=[s]{/*}{*/},
	morestring=[b]",
}


\lstdefinelanguage{SpecElektra}{
	%
	comment=[l]{;},
	commentstyle=\color{purple}\ttfamily,
	%
	morestring=[b]',
	morestring=[b]`,
	morestring=[b]",
	stringstyle=\color{purple}\ttfamily,
	%
	sensitive=f,% keywords are not case sensitive
	%
	% Colors see https://en.wikibooks.org/wiki/LaTeX/Colors
	%
	keywordstyle=\color{BlueViolet}\bfseries,
	keywordstyle=[2]\color{Green},
	keywordstyle=[3]\color{Aquamarine}\bfseries\textit,
	keywordstyle=[4]\color{NavyBlue}\bfseries,
	keywordstyle=[5]\color{Mahogany},
	%
	keywords={layer, require, validation, check, range, description, rationale, requirement, accessibility, enable, condition, message, default, opt, readonly, type, context, property1, property2, description, file, content, mountpoint, metadata, infos, plugins},
	keywords=[2]{},
	keywords=[3]{order, interface, network, emphasized},
	%keywords=[4]{[, ]},  %Not needed
	keywords=[4]{},
	keywords=[5]{},
	%
	literate={:=}{{{\color{red}\textbf:=}}}2
		 {\%}{{{\color{NavyBlue}\textbf\%}}}1
		 {[}{{{\color{Sepia}\textbf[}}}1
		 {]}{{{\color{Sepia}\textbf]}}}1,
}

\lstdefinelanguage{Cpp}{%
	language     = C++,
	literate=
}


\lstdefinelanguage{CfgElektra}{
	comment=[l]{;},
	commentstyle=\color{purple}\ttfamily,
	%
	morestring=[b]',
	morestring=[b]`,
	morestring=[b]",
	stringstyle=\color{purple}\ttfamily,
	%
	%
	sensitive=f,% keywords are not case sensitive
	%
	% Colors see https://en.wikibooks.org/wiki/LaTeX/Colors
	%
	keywordstyle=\color{Bittersweet}\bfseries,
	keywordstyle=[2]\color{DarkOrchid}\bfseries,
	keywordstyle=[3]\color{ForestGreen}\bfseries\textit,
	keywordstyle=[4]\color{Goldenrod}\bfseries,
	keywordstyle=[5]\color{CarnationPink},
	%
	keywords={},
	keywords=[2]{},
	keywords=[3]{},
	keywords=[4]{},
	keywords=[5]{},
	%
	literate={=}{{{\color{ForestGreen}\textbf=}}}1
		 %{<-}{{{\color{ForestGreen}\textbf<-}}}2
		 %{*}{{{\color{Bittersweet}\textbf*}}}1
		 {\%}{{{\color{NavyBlue}\textbf\%}}}1,
}




\lstset{language=SpecElektra, % Use SpecElektra as default programming language
	%boxpos=t, % make boxes a bit more unbreakable
	%frame=lines, % top+bottom line
	basicstyle=\ttfamily, % Use normal-size true type font
	showspaces,%
	showstringspaces=false, % Don't put marks in string spaces
	showlines=true, % make sure empty lines at end are shown (needed for concurrency
	tabsize=4, % spaces per tab
	xleftmargin=\parindent, % should be 18pt or 1.5em as defined by memoir
	%Does not really work well (needs to be deactivated for shortlistings):
	breaklines=false,
	%postbreak=\mbox{\textcolor{red}{$\hookrightarrow$}\space},
	%breakautoindent=true,
	%prebreak={\mbox{\ensuremath{\curvearrowright}}} % Zeichen am Zeilenende (Umbruch)
	%breaklines=true,
	%breakautoindent=true,
	%prebreak=\small\symbol{'134}, % backslash
	%prebreak={\mbox{\ensuremath{\curvearrowright}}} % lange kure
	%prebreak={\mbox{\ensuremath{\hookleftarrow}}} % lange kure
	%xleftmargin=3.0ex, %for some formats
	%xrightmargin=1.0ex, %for some formats
	%
	% Files do not work in utf8 see also:
	% http://stackoverflow.com/questions/1116266/listings-in-latex-with-utf-8-or-at-least-german-umlauts
	% http://tex.stackexchange.com/questions/24528/having-problems-with-listings-and-utf-8-can-it-be-fixed
	% Should work but doesn't? (Maybe add to literate broken?)
	%add to literate={ö}{{\"o}}1
	%	{ä}{{\"a}}1
	%	{ü}{{\"u}}1
	%	{Ö}{{\"O}}1
	%	{Ä}{{\"A}}1
	%	{Ü}{{\"U}}1
	%	{ß}{{\ss}}1,
	%
	% listingsutf8 did not work, made umlauts in comments very strange
	%extendedchars=true,
	%inputencoding=utf8,
	%
	%morecomment=[l][\color{blue}]{...}, % Line continuation (...) e.g. blue comment
	morekeywords={for_each},
	numbers=left, % Line numbers on left
	firstnumber=1, % Line numbers start with line 1
	numberstyle=\small\color{blue}, % Line numbers are blue and small
	numbersep=5pt,
	%stepnumber=5 % Line numbers go in steps of 5
}



\lstMakeShortInline[postbreak=,keywordstyle={}]^

\graphicspath{{../pic/}{../figures/}{../graphics/}{../ipe/}{../ggplot/}}



\lstdefinelanguage{dump}
{
	morekeywords={kdbOpen,ksNew,keyNew,keyMeta,keyCopyMeta,keyEnd,ksEnd,kdbClose},
	sensitive=false,
	morecomment=[l]{//},
	morecomment=[s]{/*}{*/},
	morestring=[b]",
}


\lstdefinelanguage{SpecElektra}{
	%
	comment=[l]{;},
	commentstyle=\color{purple}\ttfamily,
	%
	morestring=[b]',
	morestring=[b]`,
	morestring=[b]",
	stringstyle=\color{purple}\ttfamily,
	%
	sensitive=f,% keywords are not case sensitive
	%
	% Colors see https://en.wikibooks.org/wiki/LaTeX/Colors
	%
	keywordstyle=\color{BlueViolet}\bfseries,
	keywordstyle=[2]\color{Green},
	keywordstyle=[3]\color{Aquamarine}\bfseries\textit,
	keywordstyle=[4]\color{NavyBlue}\bfseries,
	keywordstyle=[5]\color{Mahogany},
	%
	keywords={layer, require, validation, check, range, description, rationale, requirement, accessibility, enable, condition, message, default, opt, readonly, type, context, property1, property2, description, file, content, mountpoint, metadata, infos, plugins},
	keywords=[2]{},
	keywords=[3]{order, interface, network, emphasized},
	%keywords=[4]{[, ]},  %Not needed
	keywords=[4]{},
	keywords=[5]{},
	%
	literate={:=}{{{\color{red}\textbf:=}}}2
		 {\%}{{{\color{NavyBlue}\textbf\%}}}1
		 {[}{{{\color{Sepia}\textbf[}}}1
		 {]}{{{\color{Sepia}\textbf]}}}1,
}

\lstdefinelanguage{Cpp}{%
	language     = C++,
	literate=
}


\lstdefinelanguage{CfgElektra}{
	comment=[l]{;},
	commentstyle=\color{purple}\ttfamily,
	%
	morestring=[b]',
	morestring=[b]`,
	morestring=[b]",
	stringstyle=\color{purple}\ttfamily,
	%
	%
	sensitive=f,% keywords are not case sensitive
	%
	% Colors see https://en.wikibooks.org/wiki/LaTeX/Colors
	%
	keywordstyle=\color{Bittersweet}\bfseries,
	keywordstyle=[2]\color{DarkOrchid}\bfseries,
	keywordstyle=[3]\color{ForestGreen}\bfseries\textit,
	keywordstyle=[4]\color{Goldenrod}\bfseries,
	keywordstyle=[5]\color{CarnationPink},
	%
	keywords={},
	keywords=[2]{},
	keywords=[3]{},
	keywords=[4]{},
	keywords=[5]{},
	%
	literate={=}{{{\color{ForestGreen}\textbf=}}}1
		 %{<-}{{{\color{ForestGreen}\textbf<-}}}2
		 %{*}{{{\color{Bittersweet}\textbf*}}}1
		 {\%}{{{\color{NavyBlue}\textbf\%}}}1,
}




\lstset{language=SpecElektra, % Use SpecElektra as default programming language
	%boxpos=t, % make boxes a bit more unbreakable
	%frame=lines, % top+bottom line
	basicstyle=\ttfamily, % Use normal-size true type font
	showspaces,%
	showstringspaces=false, % Don't put marks in string spaces
	showlines=true, % make sure empty lines at end are shown (needed for concurrency
	tabsize=4, % spaces per tab
	xleftmargin=\parindent, % should be 18pt or 1.5em as defined by memoir
	%Does not really work well (needs to be deactivated for shortlistings):
	breaklines=false,
	%postbreak=\mbox{\textcolor{red}{$\hookrightarrow$}\space},
	%breakautoindent=true,
	%prebreak={\mbox{\ensuremath{\curvearrowright}}} % Zeichen am Zeilenende (Umbruch)
	%breaklines=true,
	%breakautoindent=true,
	%prebreak=\small\symbol{'134}, % backslash
	%prebreak={\mbox{\ensuremath{\curvearrowright}}} % lange kure
	%prebreak={\mbox{\ensuremath{\hookleftarrow}}} % lange kure
	%xleftmargin=3.0ex, %for some formats
	%xrightmargin=1.0ex, %for some formats
	%
	% Files do not work in utf8 see also:
	% http://stackoverflow.com/questions/1116266/listings-in-latex-with-utf-8-or-at-least-german-umlauts
	% http://tex.stackexchange.com/questions/24528/having-problems-with-listings-and-utf-8-can-it-be-fixed
	% Should work but doesn't? (Maybe add to literate broken?)
	%add to literate={ö}{{\"o}}1
	%	{ä}{{\"a}}1
	%	{ü}{{\"u}}1
	%	{Ö}{{\"O}}1
	%	{Ä}{{\"A}}1
	%	{Ü}{{\"U}}1
	%	{ß}{{\ss}}1,
	%
	% listingsutf8 did not work, made umlauts in comments very strange
	%extendedchars=true,
	%inputencoding=utf8,
	%
	%morecomment=[l][\color{blue}]{...}, % Line continuation (...) e.g. blue comment
	morekeywords={for_each},
	numbers=left, % Line numbers on left
	firstnumber=1, % Line numbers start with line 1
	numberstyle=\small\color{blue}, % Line numbers are blue and small
	numbersep=5pt,
	%stepnumber=5 % Line numbers go in steps of 5
}



\lstMakeShortInline[postbreak=,keywordstyle={}]^

\graphicspath{{../pic/}{../figures/}{../graphics/}{../ipe/}{../ggplot/}}




\date{13.03.2019}

\begin{document}

\renewcommand{\enquote}[1]{\emph{``#1''}} % Cannot be done earlier

%%%%%%%%%%%%%%%%%%%%%%%%%%%%%%%
\begin{frame}
	\titlepage
	\doclicenseThis
\end{frame}

%%%%%%%%%%%%%%%%%%%%%%%%%%%%%%%%%%%%%%%%%% 
\section{Background}
{
\shadowcolor{black}
\shadowoffset{0.5pt}
\usebackgroundtemplate{\includegraphics[width=\paperwidth]{pics/clouds.jpg}}%
\begin{frame}
	\frametitle{\shadowtext{Misconfiguration}}
	\begin{itemize}[<+->]
		\item   \textcolor{white}{\shadowtext{\empha[misconfiguration]{misconfigurations}~\cite{yin2011empirical,su2007autobash,attariyan2010automating,xu2015systems} are a major cause}}\\
			\textcolor{white}{\shadowtext{of system failures~\cite{wool2004quantitative,oppenheimer2003internet,pertet2005causes}}}
		\item   \textcolor{white}{\shadowtext{much time is needed to fix misconfigurations~\cite{rabkin2011static,oppenheimer2003internet,yin2011empirical,mahajan2002bgp}}}
	\end{itemize}
\end{frame}
}

\begin{frame}
	\frametitle{Elektra~\cite{raab2016elektra}}
	\begin{itemize}[<+->]
	%	\item \citet{holland2001nofutz}~defined \empha{futzing} to denote \enquote{tinkering or fiddling experimentally}
	%	\item With \intro[no-futz computing]{no-futz computing} \citet{holland2001nofutz} mean \enquote{that futzing should be allowed, but should never be required}
		\item \intro{configuration specification languages} mitigate misconfigurations
		\item \elektra{} is a framework implementing a modular \intro{configuration specification language} for configuration settings
	%	\item In \elektra{} we use properties to specify configuration settings and configuration access
		\item \elektra{} enables \intro[no-futz computing]{no-futz computing} \cite{holland2001nofutz},
			i.e., error-prone \enquote{tinkering or fiddling experimentally} \enquote{should be allowed, but should never be required}
	\end{itemize}
\end{frame}

\begin{frame}
	\frametitle{Requirements of Elektra~\cite{raab2017challenges}}

	\begin{itemize}
		\item extensible and flexible
		\item external to applications
		\item open for introspection
	\end{itemize}
\end{frame}

\section{Elektra}

\begin{frame}
	\frametitle{Metalevels}
	\includegraphics{metalevels}

	We will now walk through metalevels bottom-up.
\end{frame}

\begin{frame}[fragile]
	\frametitle{Configuration Settings}

	A configuration file may look like (properties format):

	\begin{code}[language=CfgElektra]
	slapd/threads/listener=4
	\end{code}

	\vspace{1cm}

	We apply these configuration settings imperatively using:

	\begin{code}[language=bash]
	kdb set /slapd/threads/listener 4
	\end{code}
\end{frame}

\begin{frame}[fragile]
	\frametitle{Specifications}
	For specifications such as:

	\begin{code}
	[slapd/threads/listener]
	  check/range:=1,2,4,8,16
	  default:=1
	  visibility:=advanced
	  description:=One thread is adequate\
		       for up to 16 CPU cores.
	\end{code}

	\vspace{0.6cm}

	We apply the specifications imperatively using:

	\begin{code}[language=bash,morekeywords={meta-set}]
	kdb meta-set /slapd/threads/listener\
		check/range 1,2,4,8,16
	kdb meta-set /slapd/threads/listener\
		default 1
	\end{code}
\end{frame}

\begin{frame}[fragile]
	\frametitle{Meta-Specifications}
	For meta-specifications such as:

	\small
	\begin{code}
	[visibility]
	type:=enum critical important user\
	      advanced developer debug disabled
	description:=Who should see this\
	     configuration setting?
	\end{code}

	\vspace{1cm}

	We apply the meta-specifications imperatively using:

	\begin{code}[language=bash,morekeywords={meta-set}]
	kdb meta-set /elektra/meta/\
		visibility type enum ...
	kdb meta-set /elektra/meta/\
		visibility description "Who ...
	\end{code}
\end{frame}


\begin{frame}[fragile]
	\frametitle{Grammar}
	\begin{alertblock}{Idea}
	Use configuration file format grammar to describe both configurations and (meta-)specifications
	\end{alertblock}

	\begin{grammar}
	<KeySet> ::= \lq ksNew'\WhiteSpace(' \{ <Key> \lq , \LineBreak'  \}  \{ \lq\WhiteSpace' \} \lq KS\_END);'

	<Key> ::= \lq keyNew \WhiteSpace ('' ' <key name> \lq ''  , \LineBreak' [ <Value> ] <properties> \lq KEY_END)'

	<Value> ::=  \{ \lq\WhiteSpace' \} \lq KEY\_VALUE, \WhiteSpace '' ' <configuration value> \lq ''  , \LineBreak'

	<properties> ::= \{ \{ \lq\WhiteSpace' \} <property> \lq , \LineBreak' \}

	<property> ::=  \lq KEY\_META, \WhiteSpace " ' <property name> \lq "  , \WhiteSpace " ' <property value> \lq " '
	\end{grammar}
\end{frame}

\begin{frame}[fragile]
	\frametitle{Example}
	\begin{example}
	Given the key ^/slapd/threads/listener^, with the configuration value ^4^ and the property $\property{default} \mapsto 1$, \elektra{} emits:

	\begin{code}[gobble=4,language=Cpp]
	ksNew (keyNew ("/slapd/threads/listener",
		       KEY_VALUE, "4",
		       KEY_META, "default", "1",
		       KEY_END),
	       KS_END);
	\end{code}
	\vspace{-1em}
	\end{example}

	\pause
	\begin{alertblock}{Finding}
	We have source code representing the settings.
	If we instantiate it, we get a data structure representing the settings.
	Plugins emitting such ``configuration files'' are code generators.
	\end{alertblock}
\end{frame}

\begin{frame}[fragile]
	\frametitle{Usage in Applications}

	With the specification:
	\par
	\begin{code}[gobble=4]
	[slapd/threads/listener]
	  check/range:=1,2,4,8,16
	  default:=1
	  visibility:=advanced
	  restrict/write:=1
	\end{code}
	\par
	\elektra{Gen} gives the user read-only access to the object ^env.slapd.threads.listener^:
	\par
	\begin{code}[language=Cpp]
	std::cout << env.slapd.threads.listener;
	env.slapd.threads.listener = 3; // error
	\end{code}
	\par
\end{frame}

\begin{frame}
	\frametitle{Implementation}
	\begin{columns}[c]
	\begin{column}{6cm}
	\includegraphics[scale=0.8]{horizontalmodularity}
	\end{column}
	\begin{column}{5cm}
	Cylinders are configuration files, P? are plugins~\cite{raab2016improving}.

	\begin{itemize}
	\item syntax is defined via plugins reading/writing configuration files
	\item semantics are defined via
		\begin{itemize}
		\item plugins interpreting properties
		\item generated code used by applications
		\end{itemize}
	\end{itemize}
	\end{column}
	\end{columns}
\end{frame}

%%%%%%%%%%%%%%%%%%%%%%%%%%%%%%%%%%%%%%%%%% 
\section{Configuration Management}

\begin{frame}[fragile]
	Key/value access in Chef:
	\vspace{0.5cm}

	\begin{code}[morekeywords={kdbset,do,action,value,end},gobble=4]
	kdbset '/slapd/threads/listener' do
		value '4'
		action :create
	end
	\end{code}

	\pause
	\begin{alertblock}{Finding}
	We have configuration management code representing the settings.
	\end{alertblock}
\end{frame}

\begin{frame}[fragile]
	Key/value access in Ansible:
	\vspace{0.5cm}

	\begin{code}[morekeywords={name,connection,key,value,elektra,mountpoint,file,plugins,hosts,tasks},gobble=4]
	- name: setup LDAP
	  connection: local
	  hosts: localhost
	  tasks:
	  - name: set listening threads
	    elektra:
	      key: '/slapd/threads/listener'
	      value: '4'
	\end{code}
\end{frame}

\begin{frame}[fragile]
	Key/value access in puppet-libelektra~\cite{raab2020unified}:
	\vspace{0.5cm}

	\begin{code}[morekeywords={kdbkey,kdbmount,ensure,value},gobble=4]
	kdbkey {'/slapd/threads/listener':
		ensure => 'present',
		value => '4'
		check => {
			'type' => 'short',
			'range' => '1,2,4,8,16',
			'default' => '1'
		}
	}
	\end{code}
\end{frame}


%%%%%%%%%%%%%%%%%%%%%%%%%%%%%%%%%%%%%%%%%%%%%%%%%%%%
\section{Problems with Flexibility}

\subsection{Feature Interaction}

\begin{frame}[fragile]
	\frametitle{Fixed Order}

	\begin{code}
	[foo/bar]
	  encrypt:=1
	  zip:=1
	\end{code}

	\begin{itemize}
	\item zip must happen before encrypt
	\item must be specified within plugins
	\end{itemize}
\end{frame}

\begin{frame}[fragile]
	\frametitle{Plugins}

	Plugins that modify the KeySet can be very problematic, e.g. notification:
	\begin{code}[gobble=4]
	[foo/bar]
	  rename/toupper:=1
	  notify:=1
	  type:=string
	\end{code}

	\begin{itemize}
	\item expected: notify on change of key
	\item actual: always notifies because of renaming
	\item ordering does not help, as notification must happen at end
	\end{itemize}
\end{frame}

\begin{frame}[fragile]
	\frametitle{Plugins}

	Or even plain contradicting, e.g. assignment:
	\begin{code}[gobble=4]
	[foo/bar]
	  assign/math:=1+1
	  assign/python:=1+2
	  type:=long
	\end{code}

	\begin{alertblock}{Open Question}
	How to deal with it?
	\end{alertblock}
\end{frame}


\subsection{Integration of Operators}

\begin{frame}[fragile]
	\frametitle{In C++/RelaxNG Style}

	\begin{code}[morekeywords={check,type,range,math},gobble=4]
	[t/a]
	type:=long
	restrict/write:=1
	[t/x]
	check/enum/#0:=a
	check/enum/#1:=b
	check/enum/#2:=c
	\end{code}

	could be rendered by a storage plugin by:

	\begin{code}[language=C++,morekeywords={@NonNull,@ReadOnly,string},gobble=4]
	t {
		@restrict/write long a;
		enum {a,b,c} x;
	};
	\end{code}
\end{frame}


\begin{frame}[fragile]
	\frametitle{Types}

	\begin{code}[morekeywords={check,type,range,math},gobble=4]
	[a]
	  check/type:=long
	[b]
	  check/type:=long
	[c]
	  check/type:=long
	  check/range:=0-10
	  check/math:="== + a b"
	\end{code}

	ideal to further restrict types
\end{frame}

\begin{frame}[fragile]
	\frametitle{Integration}

	\begin{code}[morekeywords={check,type,range,math},gobble=4]
	[c]
	  check/type:=long
	  check/range:=0-10
	  check/math:="== + a b"
	  check/condition "(a == '3') ? (b == '5')"
	\end{code}

	\begin{itemize}
	\item does not allow math expression to be embedded in condition
	\item temporary variables are needed
	\end{itemize}

	\begin{alertblock}{Open Question}
	How to deal with it?
	\end{alertblock}
\end{frame}


\section{Conclusion}

\begin{frame}
	\frametitle{Conclusion}
	\begin{itemize}
	\item Elektra makes configuration management tasks simpler
	\item syntax flexible but limited to key/value formats
	\item plugin development needs careful consideration of interactions
	\item syntax of operators not always combinable
	\end{itemize}
\end{frame}



%%%%%%%%%%%%%%%%%%%%%%%%%%%%%%%%%%%%%%%%%% 
\nocite{raab2017introducing}

\appendix

\begin{frame}[allowframebreaks]
	\bibliographystyle{plainnat}
	\bibliography{../shared/elektra.bib}
\end{frame}

\end{document}


